\chapter{Same-carrier data and Phase 8 consumers}
\label{route-ch-handoff}

\section{The public data boundary}

\begin{definition}[Exact \(\operatorname{PicRepDatum}\) and
    \(\operatorname{JacobianData}\) carriers]
  \label{route-def-public-data}
  \lean{AlgebraicGeometry.PicRepDatum,
    AlgebraicGeometry.PicRepDatum.grpObj,
    AlgebraicGeometry.PicRepDatum.homEquiv,
    AlgebraicGeometry.JacobianData,
    AlgebraicGeometry.JacobianData.grpObj,
    AlgebraicGeometry.JacobianData.homEquiv,
    AlgebraicGeometry.JacobianData.uniqueUpToIso}
  \leanok
  \dcref{custom}

  For fields \(k\to k'\) and a curve \(C'/k'\), a
  \(\operatorname{PicRepDatum}(k,k',C')\) consists exactly of
  \[
    J'/k',\qquad
    r':h_{J'}\xrightarrow{\sim}\operatorname{Pic}^0_{C'},\qquad
    J'\to\Spec k'\text{ locally of finite type}.
  \]
  For a curve \(C/k\), a \(\operatorname{JacobianData}(C)\) consists exactly
  of the same-field carrier and representation, the same local-finite-type
  certificate, and a quasi-compactness certificate for its structure map.
  In both structures the group object, the bijections
  \(\Hom(T,J')\simeq\operatorname{Pic}^0_{C'}(T)\), and uniqueness of the
  representing object are derived from \(r'\); they are not new fields.
\end{definition}

\begin{proof}
  This is a direct expansion of the two Lean structures.  In particular,
  neither stores properness, smoothness, a dimension, an Abel source, or an
  Albanese theorem.  Those are downstream consumers with separate inputs.
\end{proof}

\begin{theorem}[Conditional original-field package and literal handoff]
  \label{route-thm-public-handoff}
  \lean{AlgebraicGeometry.picRepDatum_finiteGaloisDescent,
    AlgebraicGeometry.jacobianData_finiteGaloisDescent,
    AlgebraicGeometry.PicRepDatum.toJacobianData,
    AlgebraicGeometry.PicRepDatum.toJacobianData_J,
    AlgebraicGeometry.PicRepDatum.toJacobianData_rep,
    AlgebraicGeometry.PicRepDatum.toJacobianData_grpObj,
    AlgebraicGeometry.PicRepDatum.homEquiv_toJacobianData}
  \leanok
  \dcref{custom}
  \uses{route-def-public-data, route-thm-finite-galois-descent}

  Under the exact hypotheses of
  Theorem~\ref{route-thm-finite-galois-descent}, let \(J/K\) be its glued
  quotient.  The locally-finite-type conclusion packages
  \((J,r)\) as \(d:\operatorname{PicRepDatum}(K,K,C)\).  Given the descended
  quasi-compactness certificate, define
  \[
    D=d.\operatorname{toJacobianData}:
       \operatorname{JacobianData}(C).
  \]
  Then, definitionally,
  \[
    D.J=d.J=J,\qquad D.\mathrm{rep}=d.\mathrm{rep}=r,
    \qquad D.\mathrm{grpObj}=d.\mathrm{grpObj},
    \qquad D.\mathrm{homEquiv}=d.\mathrm{homEquiv}.
  \]
\end{theorem}

\begin{proof}
  The checked finite-Galois wrapper builds the quotient and its representation,
  and descends local finite type and quasi-compactness from \(J_L\).  The
  `toJacobianData` constructor assigns the four displayed fields without
  transport; all four comparison lemmas are proved by reflexivity.  This is
  the required same-carrier handoff.  It remains conditional until the first
  unresolved producer in Chapter~\ref{route-ch-universal} and the downstream
  projectivity/orbit producer in Chapter~\ref{route-ch-galois} are supplied.
\end{proof}

\section{Checked consumers of an arbitrary datum}

\begin{theorem}[Universal-element, Abel, and functorial consumers]
  \label{route-thm-datum-core-consumers}
  \lean{AlgebraicGeometry.JacobianData.homEquiv,
    AlgebraicGeometry.JacobianData.homEquiv_comp,
    AlgebraicGeometry.JacobianData.ofCurve,
    AlgebraicGeometry.JacobianData.comp_ofCurve,
    AlgebraicGeometry.JacobianData.pullbackHom,
    AlgebraicGeometry.JacobianData.pullbackHom_id,
    AlgebraicGeometry.JacobianData.pullbackHom_comp,
    AlgebraicGeometry.JacobianData.baseChange,
    AlgebraicGeometry.baseChangeIsoOfData}
  \leanok
  \source{abelian-varieties:page-0110}
  \dcref{Milne III, Proposition 6.1, only for the classical pointed model}
  \uses{route-thm-public-handoff}

  For any \(D:\operatorname{JacobianData}(C)\), the representation gives a
  commutative group object on \(D.J\).  For every \(K\)-point \(P\) of \(C\),
  the class of \(\mathcal O(\Delta-P\times C)\) defines
  \[
    \operatorname{ofCurve}_D(P):C\longrightarrow D.J,
    \qquad \operatorname{ofCurve}_D(P)\circ P=e_{D.J}.
  \]
  A morphism of curves \(X\to Y\), together with data \(D_X,D_Y\), induces
  the contravariant group morphism \(D_Y.J\to D_X.J\), with checked identity
  and composition laws.  For every field extension \(K\to L\),
  \(D.\operatorname{baseChange}(L)\) has literal carrier \(D.J\times_KL\),
  and for any datum \(D_L\) for \(C_L\), uniqueness of representers gives the
  group-scheme isomorphism
  \[
    \operatorname{baseChangeIsoOfData}(D,D_L):
       D.J\times_KL\xrightarrow{\sim}D_L.J.
  \]
\end{theorem}

\begin{proof}
  All assertions are formal consequences of the represented degree-zero
  functor.  The Abel map is the inverse representing equivalence applied to
  the diagonal class; naturality proves the pointing equation.  Pullback of
  Picard classes supplies the contravariant curve map.  Field base change uses
  the Picard base-change equivalence, and the displayed isomorphism is
  uniqueness of representing objects upgraded to group objects.  Milne's
  Proposition 6.1 explains the intended pointed Jacobian map, but its Albanese
  factorization is not asserted in this theorem.
\end{proof}

\begin{theorem}[Tangent-space interface at the identity]
  \label{route-thm-tangent-interface}
  \lean{AlgebraicGeometry.JacobianData.identitySection,
    AlgebraicGeometry.JacobianData.tangentSpaceCotangentDual,
    AlgebraicGeometry.JacobianData.tangentSpaceEquivPic0Kernel,
    AlgebraicGeometry.JacobianData.cotangentSpaceDual_equiv_pic0Kernel,
    AlgebraicGeometry.JacobianData.finrank_eq_of_addEquiv_of_bijective_smul}
  \leanok
  \source{kleiman-picard}
  \dcref{Theorem 5.11}
  \uses{route-thm-datum-core-consumers}

  Let \(e:\Spec K\to D.J\) be the group identity, let
  \(R=\mathcal O_{D.J,e}\), \(\kappa(e)=R/\mathfrak m_e\), and
  \(V=(\mathfrak m_e/\mathfrak m_e^2)^\vee\).  There are checked bijections
  \[
    V\simeq T_eD.J\simeq
    \ker\!\left(
      \operatorname{Pic}^0_C(\Spec K[\epsilon])\longrightarrow
      \operatorname{Pic}^0_C(\Spec K)\right),
  \]
  These are bijections of sets.  In particular, the displayed kernel
  equivalence alone does not identify scalar multiplication or prove a
  dimension formula.  Finite-dimensionality of the cotangent space is used
  internally by the checked numerical consumers below; it is not part of this
  theorem-family claim.
\end{theorem}

\begin{theorem}[Conditional tangent dimension consumers]
  \label{route-thm-tangent-conditional}
  \lean{AlgebraicGeometry.JacobianData.finrank_cotangentSpaceDual_eq_genus,
    AlgebraicGeometry.JacobianData.nonempty_cotangentSpace_addEquiv_h1}
  \leanok
  \source{kleiman-picard}
  \dcref{Theorem 5.11 and Corollary 5.13}
  \uses{route-thm-tangent-interface}

  If
  \[
    \dim_{\kappa(e)}(\mathfrak m_e/\mathfrak m_e^2)=g(C),
    \tag{TD}
  \]
  then \(\dim_{\kappa(e)}V=g(C)\).  If in addition
  \(H^1(C,\mathcal O_C)\) is finite-dimensional over \(K\), there exists an
  additive equivalence
  \[
    \mathfrak m_e/\mathfrak m_e^2\simeq_+
      H^1(C,\mathcal O_C).
  \]
  Both checked theorems take (TD) as an explicit hypothesis.
\end{theorem}

\begin{theorem}[Semilinear epsilon-kernel comparison]
  \label{route-thm-tangent-semilinear}
  \notready
  \source{kleiman-picard}
  \dcref{Theorem 5.11, specialized to the degree-zero etale Picard functor}
  \uses{route-thm-tangent-interface}

  Construct a bijection \(i:\kappa(e)\to K\), an additive equivalence
  \[
    j:V\xrightarrow{\sim_+}H^1(C,\mathcal O_C),
  \]
  and prove, for every \(a\in\kappa(e)\) and \(v\in V\),
  \[
    j(a v)=i(a)j(v).
    \tag{SL}
  \]
  Then the checked scalar-bijection finrank lemma gives (TD), which feeds the
  two conditional consumers above.
\end{theorem}

\begin{proof}
  \emph{Proof obligation.}  Kleiman Theorem 5.11 identifies
  the tangent space of the Picard scheme with \(H^1(C,\mathcal O_C)\) by
  analysing line bundles over the dual numbers.  The AJCR specialization must
  perform that calculation inside the epsilon-kernel of the exact
  etale-sheafified degree-zero functor, prove additivity, and identify the
  scalar action through the rational identity section.  The present Lean
  equivalence is set-valued, so it cannot supply (SL) or (TD).
\end{proof}

This tangent branch is the selected numerical input for the Phase 8 dimension
producer below.  It is not an alternate construction of the Jacobian: it starts
from the same original-field datum and supplies only the missing relative
dimension count.

\begin{theorem}[Conditional Abel map under field base change]
  \label{route-thm-abel-basechange-conditional}
  \lean{AlgebraicGeometry.abelCrossBaseCechCore,
    AlgebraicGeometry.abelCrossBaseCechCore_of_graph,
    AlgebraicGeometry.baseChange_ofCurve_data_of_cech}
  \leanok
  \dcref{custom: Cech representatives of the diagonal and constant graph}
  \uses{route-thm-datum-core-consumers}

  For a field extension \(K\to L\), a point \(P:\Spec K\to C\), and data
  \(D,D_L\) for \(C,C_L\), assume
  \texttt{abelCrossBaseCechCore K L C P}: transport of the Cech class
  \([\mathcal O(\Delta)]-[\mathcal O(P\times C)]\) along the pullback counit
  and the scalar-tower comparison equals the corresponding class on \(C_L\).
  Then
  \[
    \operatorname{baseChangeIsoOfData}(D,D_L)\circ
      (\operatorname{ofCurve}_D(P))_L
      =\operatorname{ofCurve}_{D_L}(P_L).
    \tag{AJ-BC}
  \]
  Equality of the transported diagonal-graph and constant-graph classes is a
  sufficient checked hypothesis for the Cech core.
\end{theorem}

\begin{theorem}[The Abel Cech base-change core]
  \label{route-thm-abel-basechange-core}
  \notready
  \dcref{custom}
  \uses{route-thm-abel-basechange-conditional}

  Prove \texttt{abelCrossBaseCechCore K L C P} for every field extension
  \(K\to L\), every challenge curve \(C/K\), and every \(K\)-point \(P\).
  This is the first unmet producer for (AJ-BC).  The graph-factor theorem is a
  reduction to two equalities, not a proof of either equality.
\end{theorem}

\begin{proof}
  \emph{Proof obligation.}  Compare the local equations of
  the diagonal and constant graph before and after the nested curve/test base
  change, prove their Cech cocycles correspond, then use multiplicativity and
  inversion to obtain the Abel class equation.  Only after this theorem may
  the conditional datum-level equality be used in the challenge assembly.
\end{proof}

\begin{theorem}[Coherent base change for the selected datum family]
  \label{route-thm-phase8-basechange-coherence}
  \notready
  \dcref{custom: exact Challenge base-change interface}
  \uses{route-thm-public-handoff, route-thm-datum-core-consumers,
    route-thm-abel-basechange-core}

  Select the datum \(D_C\) from the same R8--R9 construction for every challenge
  curve and every field extension.  For \(K\to L\), define
  \[
    \gamma_{C,L}:=
      \operatorname{baseChangeIsoOfData}(D_C,D_{C_L}):
      D_C.J\times_KL\xrightarrow{\sim}D_{C_L}.J.
  \]
  Prove that \(\gamma_{C,K}\) is the canonical identity comparison and, for
  \(K\to L\to M\), that \(\gamma_{C,M}\) equals the composite of the canonical
  pullback associator with \((\gamma_{C,L})_M\) and \(\gamma_{C_L,M}\).  Prove
  these as group-scheme equalities and prove (AJ-BC) for every \(P\in C(K)\).
  These are exactly the producers for \texttt{Jacobian.baseChangeIso},
  \texttt{baseChangeIso\_id}, \texttt{baseChangeIso\_comp}, and
  \texttt{baseChange\_ofCurve}.
\end{theorem}

\begin{proof}
  The object isomorphism exists by uniqueness of representers, but no current
  datum theorem supplies its identity and tower equations.  Normalize all
  three isomorphisms by the universal Picard class and use extensionality of
  representing morphisms to prove the two coherences.  The Abel equation is
  the checked conditional theorem instantiated by the Cech producer from
  Theorem~\ref{route-thm-abel-basechange-core}.  All equations must use the
  same selected data; arbitrary isomorphic replacement is not allowed.
\end{proof}

\begin{theorem}[Exact producers for the Phase 8 geometry interface]
  \label{route-thm-phase8-geometry-producers}
  \notready
  \source{kleiman-picard}
  \dcref{Section 3, Theorems 3.7 and 3.13; Section 5, Lemma 5.1,
    Theorem 5.11, and Corollary 5.13}
  \uses{route-thm-public-handoff, route-thm-tangent-semilinear}

  For the selected original-field datum \(D_C\), construct these three exact
  inputs on its unchanged carrier.
  \begin{enumerate}
  \item Put \(n_C=\operatorname{divRepAffAdmissibleParameter}(C)\).  Construct
    the proper degree-\(n_C\) effective-divisor scheme
    \(D_C^{\mathrm{prop}}\), its geometric-irreducibility certificate, and the
    translated Abel morphism
    \[
      a_C:D_C^{\mathrm{prop}}\longrightarrow D_C.J
    \]
    obtained from the universal divisor class using the fixed
    \texttt{admissibleAbelChartIndex}/\texttt{admissibleAbelChartDivisor}
    translation.  Prove \(a_C\) surjective and package precisely
    \[
      A_C:\operatorname{AbelSourceData}(D_C)
    \]
    with fields \((D_C^{\mathrm{prop}},\operatorname{IsProper},
    \operatorname{GeometricallyIrreducible},a_C,
    \operatorname{Surjective}(a_C))\).
  \item Construct
    \(h_{\mathrm{red}}:
      \operatorname{GeometricallyReduced}(D_C.J\to\Spec K)\).
  \item Construct a Zariski zero-hypercover
    \(\mathcal U_C\) of \(D_C.J\) and, for every leg \(i\), a certificate
    \[
      \operatorname{SmoothOfRelativeDimension}
        \bigl(g(C),\mathcal U_C.f(i)\mathbin{\gg}D_C.J.\mathrm{hom}\bigr).
      \tag{DimCov}
    \]
  \end{enumerate}
\end{theorem}

\begin{proof}
  These are planned downstream producers, not assumptions to erase.  For the first,
  Kleiman's divisor/Hilbert construction is the classical source shape; the
  AJCR proof must build the actual proper divisor scheme rather than reuse the
  widened affine-test representer, define the displayed fixed translation, and
  specialize the Riemann inequality to prove surjectivity on scheme points.
  This downstream source is not an alternate construction of \(D_C.J\).

  For the second, specialize Kleiman Lemma 5.1 to the exact degree-zero functor:
  prove reducedness at the identity from the dual-number computation and move
  it by the checked translations, including arbitrary residue-field points.
  For the third, use (SL) to fix the identity-chart dimension at \(g(C)\), use
  geometric reducedness and the group-scheme smoothness criterion to obtain a
  standard-smooth identity chart, and construct the translated Zariski cover
  with (DimCov).  The Lean handoff is exactly the three arguments consumed
  below; no theorem currently constructs them.
\end{proof}

\begin{theorem}[Conditional geometric consumers]
  \label{route-thm-datum-geometric-consumers}
  \lean{AlgebraicGeometry.JacobianData.isSeparated,
    AlgebraicGeometry.AbelSourceData,
    AlgebraicGeometry.isProper_of_abelSource,
    AlgebraicGeometry.geometricallyIrreducible_of_abelSource,
    AlgebraicGeometry.JacobianData.smooth,
    AlgebraicGeometry.JacobianData.isAbelianVariety_of_abelSource,
    AlgebraicGeometry.JacobianData.smoothOfRelativeDimension_of_translation_cover}
  \leanok
  \source{kleiman-picard}
  \dcref{Section 5, Lemma 5.1 and Corollaries 5.13--5.14}
  \uses{route-thm-phase8-geometry-producers}

  Every datum carrier is separated.  If an
  \(\operatorname{AbelSourceData}(D)\) is supplied---a proper geometrically
  irreducible \(K\)-scheme mapping surjectively to \(D.J\)---then \(D.J\) is
  proper and geometrically irreducible.  If \(D.J\) is geometrically reduced,
  then it is smooth.  A Zariski zero-hypercover whose legs are smooth of
  relative dimension \(g(C)\) proves
  \[
    \operatorname{SmoothOfRelativeDimension}(g(C),D.J\to\Spec K).
  \]
  The name \texttt{smoothOfRelativeDimension\_statement} is only an
  abbreviation for this proposition, not a proof or a producer.
\end{theorem}

\begin{proof}
  The unit of a group scheme over a field is closed, so the diagonal is closed
  and the carrier is separated.  Properness and geometric irreducibility
  descend along the surjective Abel source in the checked consumer theorems.
  A geometrically reduced group scheme locally of finite type over a field is
  smooth.  Relative dimension is local on the source, so it glues from the
  displayed hypercover.  Each implication is checked, but none constructs its
  own Abel source, geometric-reducedness proof, or dimension chart.
\end{proof}

\section{The blocked Phase 8 assembly}

\begin{theorem}[Rational-extension and uniqueness substrate]
  \label{route-thm-albanese-substrate}
  \lean{AlgebraicGeometry.Scheme.RationalMap.indeterminacy_pure_codim_one_into_grpScheme,
    AlgebraicGeometry.Scheme.RationalMap.extend_to_av,
    AlgebraicGeometry.JacobianData.existsUnique_ofCurve_comp_of_baseChange,
    AlgebraicGeometry.JacobianData.existsUnique_ofCurve_comp_of_baseChange_datum}
  \leanok
  \source{abelian-varieties:page-0023, abelian-varieties:page-0024}
  \dcref{Milne I, Theorem 3.2 and Lemma 3.3}
  \uses{route-thm-datum-core-consumers}

  Over an algebraically closed field, a rational map from a smooth integral
  variety to an abelian variety extends uniquely once the displayed
  pure-codimension-one indeterminacy input is supplied; the checked Milne--3.3
  theorem supplies that input under its stated variety and group hypotheses.
  Separately, if
  \[
    \exists u:D.J\to A,\quad f=u\circ\operatorname{ofCurve}_D(P)
    \tag{Hex}
  \]
  and uniqueness of the corresponding morphism is known after base change,
  the checked faithful-base-change theorem upgrades (Hex) to the required
  \(\exists!\) over \(K\).  It also accepts an independently constructed datum
  over the extension, using uniqueness of representers to compare carriers.
\end{theorem}

\begin{proof}
  Milne Theorem I.3.2 combines codimension at least two for maps into proper
  targets with Lemma I.3.3, which says that nonempty indeterminacy into a group
  variety is pure of codimension one; hence the locus is empty and the rational
  map is regular.  For the descent statement, flat surjective base change is
  faithful on morphisms, so two factors equal after base change are equal over
  \(K\).  Neither argument constructs the factor in (Hex); that is a separate
  positive-genus obligation.
\end{proof}

\begin{theorem}[Conditional genus-zero package]
  \label{route-thm-genus-zero-conditional}
  \lean{AlgebraicGeometry.JacobianData.isTerminal_of_pic0Subgroup_eq_bot,
    AlgebraicGeometry.jacobianData_of_vanishing,
    AlgebraicGeometry.isTerminal_jacobianData_of_vanishing,
    AlgebraicGeometry.JacobianData.existsUnique_ofCurve_comp_of_pic0Subgroup_eq_bot,
    AlgebraicGeometry.existsUnique_ofCurve_comp_of_vanishing}
  \leanok
  \dcref{Milne I, Proposition 3.9; III, Section 2, genus-zero convention}
  \uses{route-thm-albanese-substrate}

  If
  \[
    \forall T/K,\qquad \operatorname{Pic}^0_C(T)=\{0\},
    \tag{V0}
  \]
  the checked constructor represents this functor by \(\Spec K\), and any
  datum representing it has terminal carrier.  If \(C\to\Spec K\) is
  surjective, \(P\in C(K)\), and the explicit existence statement (Hex) is
  supplied, then the genus-zero factor is unique.
\end{theorem}

\begin{proof}
  Under (V0), the Yoneda functor of the terminal object and the degree-zero
  Picard functor are both singleton-valued, giving the datum and terminality.
  The Abel map to a terminal carrier is an epimorphism for a nonempty curve, so
  precomposition proves uniqueness.  Existence is not formal: it requires the
  map \(C\to A\) to be constant.  Milne I.3, Proposition 3.9 proves this only
  for rational maps from \(\mathbb A^1\) or \(\mathbb P^1\).  Its
  specialization here therefore first requires the planned pointed
  genus-zero classification \(C\simeq\mathbb P^1_K\), or an equivalent
  rational-parametrization bridge.  No checked theorem currently proves that
  bridge, derives (V0) from \(g(C)=0\), or supplies (Hex); all three remain
  explicit inputs.
\end{proof}

\begin{theorem}[The pointed Albanese consumer]
  \label{route-thm-phase8-albanese}
  \notready
  \source{abelian-varieties:page-0110}
  \dcref{Milne III, Proposition 6.1}
  \uses{route-thm-albanese-substrate, route-thm-genus-zero-conditional,
    route-thm-phase8-geometry-producers,
    route-thm-datum-geometric-consumers}

  For a \(K\)-point \(P:\Spec K\to C\), an abelian variety \(A/K\), and
  \(f:C\to A\) with \(f\circ P=e_A\), prove
  \[
    \exists!\,u:D.J\to A,
       \qquad f=u\circ\operatorname{ofCurve}_D(P).
    \tag{Alb}
  \]
  The final Lean target is
  \texttt{Jacobian.exists\_unique\_ofCurve\_comp} after \(D.J\) is selected as
  \texttt{Jacobian C}.
\end{theorem}

\begin{proof}
  \emph{Proof obligation.}  Milne Proposition 6.1 constructs
  the factor through the degree-\(g\) symmetric sum, extends the resulting
  rational map to the abelian target, proves it is a homomorphism by rigidity,
  and uses generation by \(g\) copies of the Abel image for uniqueness.  The
  cited proposition assumes \(g>0\); the challenge theorem does not.  The
  genus-zero terminal case and the positive-genus construction must therefore
  be joined explicitly.  Existing rational-extension, rigidity, and
  conditional uniqueness files are consumers, not a theorem of type (Alb).
\end{proof}

\begin{theorem}[Phase 8 geometric and base-change assembly]
  \label{route-thm-phase8-assembly}
  \notready
  \source{kleiman-picard, abelian-varieties:page-0110}
  \dcref{Kleiman Section 5; Milne III, Section 6}
  \uses{route-thm-public-handoff, route-thm-datum-core-consumers,
    route-thm-tangent-semilinear, route-thm-phase8-basechange-coherence,
    route-thm-phase8-geometry-producers,
    route-thm-datum-geometric-consumers, route-thm-phase8-albanese}

  Select one original-field datum \(D_C\) for every challenge curve, with
  carrier definitionally equal to \(\operatorname{Jacobian}(C)\).  Discharge,
  without changing that carrier:
  \begin{enumerate}
  \item the group object and smooth-relative-dimension-\(g(C)\) instance;
  \item properness and geometric irreducibility;
  \item \(\operatorname{ofCurve}\), its pointing law, and (Alb);
  \item the contravariant functor on \(\operatorname{Curve}(K)\);
  \item for every \(K\to L\), the group-scheme base-change isomorphism, its
    identity and tower coherences, and compatibility with
    \(\operatorname{ofCurve}\).
  \end{enumerate}
\end{theorem}

\begin{proof}
  \emph{Proof obligation.}  Instantiate the datum-level
  consumers above with the same selected \(D_C\), supply the missing geometric
  producers, and define every headline object by projection from that datum.
  Uniqueness of representers gives the base-change object isomorphism; the
  existing datum-level lemmas supply its group structure and Abel compatibility,
  after which uniqueness proves the identity and tower equations.

  None of these headline targets currently has a public proof.  The exact
  declaration-level audit, including the axiom-contaminated derived
  \texttt{Jacobian.congr}, is recorded in
  Appendix~\ref{route-app-challenge-status}; it supplies no mathematical
  premise to this theorem.
\end{proof}
